\section{$q$-Dependence Audit of CoreFalcon+ Security Bound}

This section re-derives the security bound for CoreFalcon+, examining term-by-term 
where exactly the modulus plays a role in their argument. The derivations in this 
audit originate from the CoreFalcon+ paper.

\paragraph{\texorpdfstring{\(r_u\): random-oracle / syndrome
uniformity}{r\_u: random-oracle / syndrome uniformity}}\label{r_u-random-oracle-syndrome-uniformity}

\(r_u\) controls the Rényi divergence between the uniform distribution
\(U(R_q)\) and the distribution of the programmed syndrome
\(s_1 + s_2 h \bmod q\) for Gaussian \(s_1, s_2 \sim D_{R,s}\):

\[
r_u = \max_{(\cdot,h)\in \mathrm{TpdGen}} R_{a_u}(P \,\|\, Q_h), \qquad P = U(R_q).
\]

Where \(Q_h\) is the distribution of \(s_1 + s_2\cdot h \bmod q\) with
\(s_1, s_2 \sim D_{R, s}\) and signature standard deviation

\[
s = \frac{1}{\pi} \sqrt{\frac{\ln\!\left(4n\left(1+\frac{1}{\epsilon}\right)\right)}{2}} \cdot \alpha \sqrt{q}.
\]

By \hyperref[cor:ru-renyi]{Corollary~\ref*{cor:ru-renyi}}~\cite{Fouque2024CoreFalcon}
(proof in Appendix B.3 of \cite{Fouque2024CoreFalcon}), with \(P = U(R_q)\) and
\(Q\) the distribution of \(u + v \cdot h \mod q\) for \(u, v \sim D_{R,s}\),
the syndrome-uniformity R\'enyi divergence \(R_a(P \,\|\, Q)\) obeys the bound
stated there. The only nontrivial pre-condition behind it is 
\(s \ge \eta_{\epsilon}(\Lambda_{h,q})\). Bases \(B_h\) and \(B_{f,g}\)
of \(\Lambda_{h, q}\) are defined by
\(TpdGen(\alpha, q) \rightarrow (f, g, F, G, h)\) where
\(h \in R_q \setminus \{0\}\) and \((f, g, F, G) \in R^4\). The trapdoor
basis satisfies \(||B_{f,g}||_{GS} \le \alpha\sqrt{q}\). Proof 3 then resolves this condition, confirming

\[
r_u \lesssim 1 + \frac{2a_u\epsilon^2}{(1-\epsilon)^2}.
\]

Once the smoothing condition holds, this bound has no remaining
\(q\)-dependence beyond the cyclotomic ring setting and \(||B_{f,g}||_{GS} \le \alpha\sqrt{q}\) for 
Proof 3 to apply.

\paragraph{\texorpdfstring{\(r_p\): sampler output versus ideal
Gaussian}{r\_p: sampler output versus ideal Gaussian}}\label{r_p-sampler-output-versus-ideal-gaussian}

\(r_p\) is the Rényi divergence between the pre-image sampler
output \(\mathrm{PreSmp}(B,s,(c,0))\) and the ideal discrete Gaussian
\(D_{\Lambda,s}\)

\[
r_p = \max_{(B,\cdot)\in \mathrm{TpdGen},\, c\in R_q} R_{a_p}\!\left( \mathrm{PreSmp}(B,s,(c,0)) \,\middle\|\, D_{\Lambda,s} \right),
\]

where \(\Lambda = \Lambda(B)_{(c,0)}\) is the coset obtained by shifting
the lattice by the syndrome \((c,0)\). By
\hyperref[cor:rp-renyi]{Corollary~\ref*{cor:rp-renyi}}~\cite{Fouque2024CoreFalcon},
for \(s \ge \eta_\epsilon(\mathbb{Z}^{2n}) \cdot ||B||_{GS}\) and arbitrary
syndrome \(c \in R_q\), the R\'enyi divergence is
\(R_a(\mathrm{PreSmp}(B,s,(c,0)) \,\|\, D_{\Lambda,s}) \lesssim 1 + 2a\epsilon^2\).
This pre-condition is demonstrated to be the case in Proof 3 for the CoreFalcon+ parameter formula 
of $s$, therefore $r_p \lesssim 1 + 2a_p\epsilon^2$ applies in our case.
Beyond this, \hyperref[cor:rp-renyi]{Corollary~\ref*{cor:rp-renyi}} makes no reference to \(q\), 
depending only on \(n\), \(\epsilon\), and \(||B||_{GS}\). Though it remains a concern whether the 
actual sampler implementation still behaves as expected with 53-bit floating point arithmetic with 
the larger 31-bit modulus. 

\paragraph{\texorpdfstring{\(p_{\mathrm{PreSmp},\beta}\): sampler
success / rejection
probability}{p\_\{\textbackslash mathrm\{PreSmp\},\textbackslash beta\}: sampler success / rejection probability}}\label{p_mathrmpresmpbeta-sampler-success-rejection-probability}

\(p_{\mathrm{PreSmp},\beta}\) is the worst-case probability of sampling
\((s_1, s_2) \in R^2\) with \(||(s_1,s_2)||_2 \le \beta\) from the FFO
sampler. The derivation below is found in appendix~E1 of
\cite{Fouque2024CoreFalcon}. We walk through it here to understand if at any point we invoke some property of $q$.
At a high level it proceeds in four steps:

\begin{enumerate}
\def\labelenumi{\arabic{enumi}.}
\tightlist
\item
  Compare \(\mathrm{PreSmp}\) to \(D_{\Lambda,s}\) using the
  relative-error bound (\hyperref[lem:ffo-relerror]{Lemma~\ref*{lem:ffo-relerror}}~\cite{Fouque2024CoreFalcon}).
\item
  Apply the Gaussian tail bound to \(D_{\Lambda+t,s}\) (\hyperref[lem:gauss-tail]{Lemma~\ref*{lem:gauss-tail}}~\cite{Fouque2024CoreFalcon}).
\item
  Use smoothing (\hyperref[lem:gauss-mass]{Lemma~\ref*{lem:gauss-mass}}~\cite{Fouque2024CoreFalcon}) to bound
  \(\rho_s(\Lambda)/\rho_s(\Lambda+t)\).
\item
  Substitute \(\beta = \tau s \sqrt{2n}\), causing the \(s\)-scale, and
  hence the \(q\)-scale, to cancel.
\end{enumerate}

Let

\[
p_{\mathrm{PreSmp},\beta}
:=
\min_{(B,\cdot)\in \mathrm{TpdGen},\, c\in R_q}
\Pr_{(s_1,s_2)\gets_{\$}\, \mathrm{PreSmp}(B,s,(c,0))}
\left[
\left\|(s_1,s_2)\right\|_2 \le \beta
\right]
\]

Take the complement

\[
\begin{aligned}
p_{\mathrm{PreSmp},\beta}
:=
\min_{(B,\cdot)\in \mathrm{TpdGen},\, c\in R_q}
1 -
\Pr_{(s_1,s_2)\gets_{\$}\, \mathrm{PreSmp}(B,s,(c,0))}
\!\left[\left\|(s_1,s_2)\right\|_2 > \beta\right]
\end{aligned}
\]

By \hyperref[lem:ffo-relerror]{Lemma~\ref*{lem:ffo-relerror}}~\cite{Fouque2024CoreFalcon},
once \(s \ge \eta_\epsilon(\mathbb{Z}^{2n}) \cdot ||B||_{GS}\) the relative error
\(\delta_{\mathrm{RE}}(\mathrm{PreSmp}(B,s,(c,0)),\, D_{\Lambda,s})\) is bounded by
approximately \(2\epsilon\).

This relative error quantifies the error in treating the discrete
Gaussian \(D_{\Lambda,s}\) as the same distribution as the distribution
of the outputs of \(\mathrm{PreSmp}(B,s,(c,0))\) over the lattice
\(\Lambda\), and is bounded by approximately \(2\epsilon\). That is to
say, we can let \(D_{\Lambda,s}\) stand in for the distribution of
\(\mathrm{PreSmp}(B,s,(c,0))\) up to a relative error bounded by
\(2\epsilon\).

\hyperref[lem:relerror-tail]{Lemma~\ref*{lem:relerror-tail}}~\cite{Fouque2024CoreFalcon}
transfers a relative-error bound \(\delta_{RE}(P,Q) = \delta\) between two
distributions on \(\mathbb{Z}^n\) into the tail bound
\(\Pr_{x \leftarrow P}[\|x\|_2 > \beta] \le \Pr_{x \leftarrow Q}[\|x\|_2 > \beta]\,(1+\delta)\)
for any \(\beta \ge 0\).

Then, by application of \hyperref[lem:ffo-relerror]{Lemma~\ref*{lem:ffo-relerror}}
to \hyperref[lem:relerror-tail]{Lemma~\ref*{lem:relerror-tail}}, we have for

\[
\delta_{\mathrm{RE}}\!\left(
    \mathrm{PreSmp}(B,s,(c,0)),\, D_{\Lambda(B)+t,s}
\right)
\lesssim
2 \epsilon
\]

\[
\Pr_{(s_1, s_2) \leftarrow \mathrm{PreSmp}(B,s,(c,0))}\!\left[\|(s_1,s_2)\|_2 > \beta\right]
\le
\Pr_{(s_1, s_2) \leftarrow D_{\Lambda(B)+t,s}}\!\left[\|(s_1,s_2)\|_2 > \beta\right] \cdot (1+2\epsilon),
\qquad
t \in \Lambda(B)_{(c,0)}.
\]

Application to our current bound gives

\[
\begin{aligned}
&\min_{(B,\cdot)\in \mathrm{TpdGen},\, c\in R_q}
1 -
\Pr_{(s_1,s_2)\gets_{\$}\, \mathrm{PreSmp}(B,s,(c,0))}
\!\left[\|(s_1,s_2)\|_2 > \beta\right]
\\
&\qquad\gtrsim
\min_{(B,\cdot)\in \mathrm{TpdGen},\, c\in R_q}
\left(
1 -
\Pr_{(s_1,s_2)\leftarrow D_{\Lambda(B)_{(c,0)},\, s}}
\!\left[\|(s_1,s_2)\|_2 > \beta\right]
\cdot (1+2\epsilon)
\right)
\end{aligned}
\]

By \hyperref[lem:gauss-tail]{Lemma~\ref*{lem:gauss-tail}}~\cite{Fouque2024CoreFalcon},
the discrete Gaussian \(D_{\Lambda+t,s}\) satisfies the tail bound
\(\Pr_{z \leftarrow D_{\Lambda + t,\, s}}[\|z\|_2 > \tau s \sqrt{n}]
\le \frac{\rho_s(\Lambda)}{\rho_s(\Lambda + t)}(\sqrt{e^{1 - \tau^2}\tau^2})^n\).

Applying to our current bound, and since the parameter
\(\beta = \tau s \sqrt{2n}\), we have

\[
\begin{aligned}
&\min_{(B,\cdot)\in \mathrm{TpdGen},\, c\in R_q}
\left(
1 -
\Pr_{(s_1,s_2)\leftarrow D_{\Lambda(B)_{(c,0)},\, s}}
\!\left[\|(s_1,s_2)\|_2 > \beta\right]
\cdot (1+2\epsilon)
\right)
\\
&\qquad\gtrsim
\min_{(B,\cdot)\in \mathrm{TpdGen},\, c\in R_q}
\left(
1 -
\frac{\rho_s(\Lambda(B))}{\rho_s\!\left(\Lambda(B)+t\right)}
\cdot
\left(\sqrt{e^{1-\tau^2}\tau^2}\right)^{2n}
\cdot
(1+2\epsilon)
\right),
\qquad t \in \Lambda(B)_{(c,0)}.
\end{aligned}
\]

By \hyperref[lem:gauss-mass]{Lemma~\ref*{lem:gauss-mass}}~\cite{Fouque2024CoreFalcon},
for \(s \ge \eta_\epsilon(\Lambda)\) the Gaussian mass satisfies
\(\rho_{s,c}(\Lambda) \in [\,1-\epsilon,\;1+\epsilon\,] \cdot \frac{(\sqrt{2\pi} \cdot s)^n}{\det(\Lambda)}\).

Since \(\Lambda(B) + t\) is a translate (coset) of \(\Lambda(B)\), the
covolume is unchanged by translation, so
\(\det(\Lambda(B)+t) = \det(\Lambda(B))\), and
\(s \ge \eta_\epsilon(\Lambda(B))\) by Proof 3. This gives

\[
p_{\mathrm{PreSmp},\beta} \ge
1 -
\left(
\frac{1+\epsilon}{1-\epsilon}
\cdot
\left(
\sqrt{e^{1-\tau^2} \tau^2}
\right)^{2n}
\cdot
(1+2\epsilon)
\right).
\]

% TODO: I think it would be good to explain the intermediate process
% of this proof that were implicit in appendix E1 to make sure it still
% applies in our context. But the algebra is the same as theirs so it
% probably can just be reduced to that proof. Probably earlier too. In
% fact, we probably can just isolate the sections of the re-derivations
% where the only remaining steps are just algebraic manipulations and can
% be used directly from existing proofs.

This lower bound depends only on \(n\), \(\tau\), and \(\epsilon\), not
on \(q\). Since \(\beta\) is always \(\tau s \sqrt{2n}\), increasing
\(q\) increases both \(s\) and \(\beta\) proportionally, and the
\(s\)-scale cancels at the \hyperref[lem:gauss-tail]{Lemma~\ref*{lem:gauss-tail}} step.

\paragraph{\texorpdfstring{Optimal Rényi orders
\(a_u, a_p, C_s\)}{Optimal Rényi orders a\_u, a\_p, C\_s}}\label{optimal-ruxe9nyi-orders-a_u-a_p-c_s}

All choices
\(a_u, a_p \in \mathbb{R}^{\ge 1}, C_s \in \mathbb{N}^{\ge 1}\) are by
definition proof constants and do not depend on \(q\).

The Rényi orders \(a_u\) and \(a_p\) are downstream of the bounds
already established for \(r_u\) and \(r_p\). Once \(r_u\) and \(r_p\)
are bounded by expressions depending only on \(\epsilon\) and the chosen
orders, optimizing the orders to minimize the overall loss introduces no
new direct dependence on \(q\).
Appendix~E2 of
\cite{Fouque2024CoreFalcon}
reports the optimal \(a_p\) as

\[
a_p =
\frac{
\lambda \epsilon^2 \ln(4)
+
\sqrt{
8C_s\lambda\epsilon^2\ln(2)
+
\lambda^2\epsilon^4\ln^2(4)
}
}{
4C_s\epsilon^2
}
\]

The optimal \(a_u\) is determined analogously from the \(r_u\) bound.
Conceptually, both \(a_u\) and \(a_p\) depend on \(\lambda\),
\(\epsilon\), \(C_s\), and the Rényi-loss optimization, not directly on
\(q\). Numerical values are reported under
\hyperref[concrete-parameter-choices]{Concrete Parameter Choices}.